\documentclass[11pt,compress,t,notes=noshow, xcolor=table]{beamer}

\input{../../style/preamble}
\input{../../latex-math/basic-math}
\input{../../latex-math/basic-ml}

\title{Optimization in Machine Learning}

\begin{document}

\titlemeta{
Univariate optimization
}{
Further remarks
}{
figure/secant.pdf
}{
\item Alternative methods
\item Strategies for non-unimodal functions
\item Software implementations
\vspace{0.5cm}
\item[] \hspace{-0.7cm}Slides based on \furtherreading{PRESS92} (Ch. 9 \& 10)
}


% ------------------------------------------------------------------------------
\begin{framei}{Alternative univariate methods}
\item GSS, quadratic interpolation, and Brent are derivative-free
\item Alternatively, if $f'$ is available: find a minimum by solving $f'(x) = 0$ using root-finding methods applied to $f'$:
\begin{itemizeM}
\item \textbf{Bisection}: halve interval, keep the half where $f'$ changes sign \\
(slow, but guaranteed convergence)
\item \textbf{Secant method}: fit a secant through the interval bounds, use its intersection with $x$-axis as $x_\text{new}$
% , keep the half where $f'$ changes sign
(fast, but may diverge)
\imageC[0.4]{figure/secant.pdf}
\item \textbf{Newton's method}: $x_{n+1} = x_n - f'(x_n)/f''(x_n)$ \\
(quadratic convergence, but requires $f''$ and fails if $f'' \approx 0$)
\end{itemizeM}
\end{framei}


% ------------------------------------------------------------------------------
\begin{framei}[sep=L]{What if $f$ is not unimodal?}
\item Throughout this chapter, we have assumed that $f$ is unimodal on the search interval
\item Without unimodality, all covered methods may fail to find the global minimum:
\begin{itemizeM}
\item \textbf{GSS / Quadratic Interpolation / Brent}: bracket shrinking discards the half with the larger function value, but the global minimum may lie in the discarded half
\item \textbf{Bisection / secant / Newton on $f'$}: find any stationary point of $f$, which may be a local minimum, local maximum, or inflection point
\end{itemizeM}
\end{framei}


% ------------------------------------------------------------------------------
\begin{framei}[sep=L]{Dealing with non-unimodal functions}
\item Practical heuristics:
\begin{itemizeM}
\item multiple restarts: run the method from several different initial brackets and return the best result
\item e.g., split the domain into sub-intervals that are approximately unimodal and apply a local method to each
\end{itemizeM}
\item Alternative: use global optimization methods (see next chapters), although finding the global minimum is still not guaranteed
\end{framei}


% ------------------------------------------------------------------------------
\begin{framei}[sep=L]{Software implementations}
\item Several standard implementations for univariate optimization implement Brent's method as default:

\begin{itemizeL}
\item Python: \texttt{scipy.optimize.minimize\_scalar()}
\item R-function \texttt{optimize()}
\item MATLAB: \texttt{fminbnd()}
\end{itemizeL}
\end{framei}


\endlecture
\end{document}